% !PRINTED 262
\weberpage{262}

%<<CONT
\weq{1}{\psi = \varphi(x_1, x_2, x_3 \ldots x_m) - \lambda x_m^2}
verschwindet. Die Determinante ist nämlich
\[
\begin{vmatrix}
a_{1,1},\, a_{1,2} \ldots a_{1,m} \\
a_{2,1},\, a_{2,2} \ldots a_{2,m} \\
\cdots\cdots\cdots\cdots \\
a_{m,1},\, a_{m,2} \ldots a_{m,m} - \lambda
\end{vmatrix}
= R_m - \lambda R_{m-1},
\]
und sie verschwindet, wenn
\[
\lambda = \frac{R_m}{R_{m-1}}
\]
gesetzt wird.

Die Function $\psi$ lässt sich dann durch $m-1$ Variable $y$ ausdrücken, und man erhält nach der Formel (6), \S.\,81
\[
\psi(x_1, x_2 \ldots x_m) = \psi(y_1, y_2 \ldots y_{m-1}, 0),
\]
oder nach (1)
\weq{2}{\varphi(x_1, x_2 \ldots x_m) = \frac{R_m}{R_{m-1}}\, x_m^2 + \varphi(y_1, y_2 \ldots y_{m-1}, 0),}
und die Untersuchung der Function $\varphi$ von $m$ Variablen ist dadurch auf die Untersuchung von $\varphi_1(y) = \varphi(y_1, y_2 \ldots y_{m-1}, 0)$ von $m-1$ Variablen zurückgeführt, deren Determinante gleich $R_{m-1}$, also von Null verschieden ist. Je nachdem $R_m : R_{m-1}$ positiv oder negativ ist, wird die Function $\varphi(x)$ ein positives oder ein negatives Quadrat mehr haben als $\varphi_1(y)$.

Durch Anwendung des gleichen Verfahrens auf $\varphi_1(y)$ und die folgenden Functionen ergiebt sich der Satz:

Die Anzahl der positiven und negativen Glieder der Reihe
\weq{3}{\frac{R_m}{R_{m-1}}, \quad \frac{R_{m-1}}{R_{m-2}}, \ldots \frac{R_1}{R_0}}
stimmt überein mit der Anzahl der positiven und negativen Quadrate, in die sich die Function $\varphi(x)$ zerlegen lässt.

Wir wollen diesem Satze noch einen etwas anderen Ausdruck geben, schicken aber folgende Erklärung voraus.

Wenn eine Reihe von Null verschiedener reeller Zahlen in bestimmter Anordnung vorliegt, so können die Vorzeichen dieser Grössen in mannigfaltiger Weise wechseln; folgen zwei Grössen von gleichem Zeichen auf einander, so findet eine \gesperrt{Zeichenfolge} (Permanenz) statt, folgt aber auf eine Grösse eine andere von entgegengesetztem Zeichen, so haben wir einen \gesperrt{Zeichenwechsel} (Variation).
